% =====================================================================
%  Registro anonimizado de la entrevista inicial (SE-01, ENT-01)
%  Fragmento LaTeX incluido desde main.tex (Anexo E).
%  No contiene nombres, telefonos, direcciones ni datos de clientes.
% =====================================================================

\subsection*{Datos de la sesi\'on}
\begin{tabularx}{\linewidth}{@{}L{4.6cm}Y@{}}
  \toprule
  Sesi\'on / instrumento & SE-01 / ENT-01 (entrevista semiestructurada) \\
  Versi\'on del guion aplicado & ENT-01 versi\'on 1.1 (Anexo~\ref{anx:entrevista}) \\
  Participante & JC-731 \\
  Rol del participante & Responsable del taller (PI-01) \\
  Fecha y hora de inicio & 21/11/2026, 14:20 \\
  Duraci\'on & Aproximadamente 7 minutos (menor que los 45--60 minutos previstos en el guion) \\
  Lugar & Local de trabajo del responsable del taller, cant\'on Quevedo \\
  Entrevistador & Basurto Anchaguano Jordan Steven \\
  Anotador & Holguin Mendoza Jonathan Alexander \\
  Consentimiento & PFC-ISR401-CONTX-A03, versi\'on 1.3, firmado el 20/11/2026 (EV-03). Autorizaciones: a) participaci\'on: s\'i; b) audio: s\'i; c) citas textuales: s\'i \\
  Grabaci\'on de audio & \pendiente{indicar si se grab\'o y si existe el archivo (EV-05)} \\
  Origen del registro & Notas tomadas durante la sesi\'on y transcritas por el equipo (EV-04). Para este registro se us\'o una versi\'on de la transcripci\'on con expresiones cotidianas en lugar de t\'erminos t\'ecnicos, sin cambiar el contenido. \pendiente{indicar si se cotej\'o con el audio} \\
  \bottomrule
\end{tabularx}

\medskip
\noindent\textit{Cada respuesta es una s\'intesis fiel de la transcripci\'on de la sesi\'on
(conservada en la carpeta privada, EV-04); no son citas literales. Los ejemplos no incluyen datos
de clientes.}

\subsection*{A. Organizaci\'on y participantes}
\begin{description}[style=nextline,leftmargin=0.6cm,font=\normalfont\bfseries]
  \item[A1. ?`Cu\'al es su funci\'on en JC\textbar{}cell?]
    Es el encargado del taller: recibe los equipos, revisa qu\'e problemas tienen, hace el
    diagn\'ostico y coordina las reparaciones.
  \item[A1 (repregunta). ?`Qu\'e decisiones toma como parte de esa funci\'on?]
    Decide por d\'onde empezar la revisi\'on, partiendo por la causa menos dif\'icil, y decide
    cu\'anto cobrar seg\'un la dificultad del trabajo y los materiales.
  \item[A2. ?`Qu\'e servicios ofrece actualmente el taller?]
    Mantenimiento de computadoras, formateo, instalaci\'on de sistemas operativos y de
    programas, limpieza interna, cambio de componentes, soluci\'on de problemas de software y de
    hardware, y revisi\'on de laptops y otros dispositivos electr\'onicos.
  \item[A2 (repregunta). ?`Desde hace cu\'anto funciona el taller?]
    Aproximadamente tres a\~nos (estimaci\'on del participante).
  \item[A3. ?`C\'omo se distribuyen las tareas entre las personas que trabajan en el taller?]
    Trabaja solo: atiende a los clientes, registra los equipos, explica los precios, coordina
    las entregas y realiza toda la parte t\'ecnica (diagn\'osticos, mantenimientos y
    reparaciones). Cuando hay bastante trabajo debe organizarse para atender todo.
  \item[A4. ?`Con qui\'enes se comunica normalmente para coordinar una reparaci\'on?]
    Con los clientes, para conocer la falla y comunicar el diagn\'ostico, y con los
    proveedores, cuando necesita conseguir un repuesto.
\end{description}

\subsection*{B. Recorrido de un servicio}
\begin{description}[style=nextline,leftmargin=0.6cm,font=\normalfont\bfseries]
  \item[B1. Reparaci\'on reciente, desde la recepci\'on hasta la entrega]
    Una laptop lenta que se apagaba de repente. Convers\'o con el cliente para entender la
    falla, revis\'o el equipo y determin\'o que necesitaba mantenimiento interno y revisi\'on
    del almacenamiento. Explic\'o lo encontrado y el costo aproximado; con la aceptaci\'on del
    cliente hizo el mantenimiento, comprob\'o el funcionamiento y avis\'o que pod\'ia retirarla.
  \item[B2. ?`Qu\'e hacen cuando reciben un equipo?]
    Lo revisa f\'isicamente y anota tipo de equipo, marca, modelo, falla y datos de contacto
    del cliente. Lo registra casi siempre en un cuaderno y algunos datos en el celular. Para
    identificar los equipos intenta mantenerlos separados y recordar a qui\'en pertenece cada
    uno; cuando se acumulan, revisa las anotaciones para no confundirlos.
  \item[B3. ?`C\'omo determinan qu\'e trabajo necesita el equipo?]
    Hace una revisi\'on inicial y pruebas. Calcula el precio seg\'un tiempo, dificultad y
    necesidad de repuestos. Comunica el diagn\'ostico y el precio de palabra o por mensaje; si
    el cliente acepta, empieza el trabajo; si no, devuelve el equipo sin intervenir.
  \item[B4. ?`C\'omo conocen el avance de cada reparaci\'on?]
    Por lo anotado en la recepci\'on y el seguimiento diario. Maneja las etapas: recibido, en
    revisi\'on, pendiente de autorizaci\'on, en reparaci\'on y terminado. No tiene un sistema
    que actualice los estados; debe revisar sus anotaciones para saber en qu\'e va cada equipo.
  \item[B5. Cuando una reparaci\'on necesita una pieza, ?`qu\'e pasos siguen?]
    Confirma la pieza, consulta disponibilidad y precio con proveedores. Si la tiene, contin\'ua;
    si no, la pide y el equipo permanece en el taller hasta que llegue.
  \item[B6. ?`C\'omo se comunican con el cliente durante el servicio?]
    Por llamadas y WhatsApp: al tener el diagn\'ostico, al necesitar autorizaci\'on para trabajo
    adicional y al terminar. Si el cliente no responde, espera y vuelve a intentar; a veces debe
    insistir porque necesita la autorizaci\'on para seguir.
  \item[B7. ?`C\'omo gestionan el cobro?]
    Cobra casi siempre al retirar el equipo, en efectivo o por transferencia. Anota en el
    cuaderno el monto y si qued\'o pagado; marca los pagos pendientes. Seg\'un el servicio
    entrega una nota o comprobante.
  \item[B8. ?`C\'omo realizan la entrega?]
    Enciende y prueba el equipo, verifica que la falla est\'e resuelta, avisa al cliente,
    confirma el pago y anota que el equipo fue retirado. Si el cliente regresa por la misma
    reparaci\'on dentro de la garant\'ia, vuelve a revisarlo.
\end{description}

\subsection*{C. Fortalezas, dificultades y evidencia}
\begin{description}[style=nextline,leftmargin=0.6cm,font=\normalfont\bfseries]
  \item[C1. ?`Qu\'e funciona bien y convendr\'ia conservar?]
    La atenci\'on directa al cliente, que le permite explicar en persona la falla, y su
    experiencia, que le ayuda a identificar fallas con rapidez.
  \item[C2. ?`En qu\'e partes encuentra dificultades?]
    Llevar el control de todos los equipos que ingresan: al registrar a mano debe buscar entre
    varias anotaciones. Le cuesta recordar qu\'e equipos est\'an pendientes de revisi\'on, cu\'ales
    esperan repuestos y cu\'ales est\'an terminados.
  \item[C3. Caso reciente]
    Un equipo qued\'o esperando un repuesto entre varios pendientes; d\'ias despu\'es tard\'o en
    recordar su estado y el trabajo faltante revisando sus notas, y perdi\'o tiempo que pod\'ia
    dedicar a otras reparaciones.
  \item[C3 (repregunta). ?`Con qu\'e frecuencia sucede?]
    No a diario; ocurre cuando hay varios equipos acumulados o una reparaci\'on se demora.
    No se indic\'o una cifra ni un registro que permita medirla.
  \item[C4. Si pudiera mejorar un solo aspecto, ?`cu\'al elegir\'ia?]
    El registro y control de los equipos: poder ver r\'apidamente qu\'e equipos tiene, de qui\'en
    son, qu\'e falla tienen y en qu\'e estado va cada reparaci\'on.
\end{description}

\subsection*{D. Resultados esperados y condiciones}
\begin{description}[style=nextline,leftmargin=0.6cm,font=\normalfont\bfseries]
  \item[D1. ?`C\'omo se dar\'ia cuenta de que la mejora funciona?]
    Ser\'ia m\'as f\'acil encontrar la informaci\'on de cada equipo y saber qu\'e reparaciones est\'an
    pendientes; habr\'ia menos confusiones y menos tiempo perdido buscando en anotaciones.
    No se mencion\'o un registro actual que permita comparar.
  \item[D2. ?`Con qu\'e recursos cuentan?]
    Celulares, una computadora y conexi\'on a Internet. La mayor parte del registro se hace a
    mano, con cuadernos y mensajes de WhatsApp.
  \item[D3. ?`Qu\'e condiciones deber\'ian respetarse?]
    Que sea sencillo, que no complique el trabajo diario y que no requiera equipos costosos ni
    conocimientos t\'ecnicos avanzados. Desea conservar la atenci\'on directa y su forma de
    diagnosticar y reparar; busca organizarse mejor, no cambiar su manera de trabajar.
  \item[D4. ?`A qui\'enes convendr\'ia consultar?]
    Considera que basta con \'el, porque realiza todas las actividades y ser\'ia el usuario
    principal de la herramienta.
\end{description}

\subsection*{E. Cierre}
\begin{description}[style=nextline,leftmargin=0.6cm,font=\normalfont\bfseries]
  \item[E1. ?`Qu\'e deber\'iamos corregir o completar del resumen?]
    El resumen refleja c\'omo trabaja; pide dejar claro que casi todos sus registros son
    manuales y que busca organizar mejor la informaci\'on.
  \item[E2. ?`Hay alg\'un tema que no hayamos tratado?]
    El seguimiento de garant\'ias y el historial de reparaciones de cada equipo.
  \item[E3. ?`C\'omo verificar la informaci\'on pendiente?]
    Revisando registros de ejemplo del taller (anotaciones de recepci\'on, trabajos realizados y
    entregas), anonimizados, y en una nueva conversaci\'on.
\end{description}

\subsection*{Pendientes de confirmar}
\begin{itemize}[nosep]
  \item Frecuencia y cantidad de casos de equipos cuyo estado debe reconstruirse (C3): no hay
        cifra ni registro; se propone revisar registros de ejemplo (E3).
  \item Los problemas mencionados en la conversaci\'on previa (p\'erdida de equipos y de
        repuestos, falta de repuestos, cobros olvidados) no fueron expresados con esos t\'erminos
        en la entrevista; deben confirmarse con el participante.
  \item Existencia de repuestos en stock y forma de registrarlos (B5 no lo aclara).
  \item Otros roles a consultar: el participante indic\'o que no hay m\'as personal (A3, D4);
        clientes y proveedores son partes interesadas externas a\'un no consultadas.
\end{itemize}
